\documentclass[10pt]{article}

\usepackage{sansmathfonts}
%\usepackage[scaled=0.86]{beramono} % for code listings
\usepackage{helvet}
\renewcommand{\rmdefault}{\sfdefault}

% \usepackage[T1]{fontenc}
% \usepackage[math]{blindtext}
% \usepackage[utf8]{inputenc}

\usepackage{amsmath}
\usepackage[dvipsnames]{xcolor}
\usepackage{xcolor-material}
\usepackage{setspace}
\usepackage{sectsty}
\usepackage{ragged2e}
\usepackage{titlesec}
\usepackage{graphicx}
\usepackage{float}
\usepackage{lastpage}
\usepackage[font={small,color=MaterialGrey500}]{caption}\usepackage{subcaption}
\usepackage{siunitx}
\usepackage{tikz}

\usepackage{geometry}
 \geometry{
 a4paper,
 left=30mm,
 right=30mm,
 top=30mm,
 bottom=30mm
 }

\captionsetup{justification=centering,singlelinecheck=false}

\usepackage{pgfplots}
\pgfplotsset{compat=newest}
\usetikzlibrary{plotmarks}
\usetikzlibrary{arrows.meta}
\usepgfplotslibrary{patchplots}
\usepackage{grffile}

\pgfplotsset{every axis/.append style={
                    label style={font=\tiny},
                    }}
                    
\AtBeginDocument{\sisetup{math-rm=\mathrm, text-rm=\rmfamily}}

\DeclareSIUnit\denier{denier}
\DeclareSIUnit\one{1}

\pgfplotsset{
    axis line style={MaterialGrey500},
    every axis label/.append style ={MaterialGrey500},
    every tick label/.append style={MaterialGrey500}
  }

\setlength\parindent{0pt}
\setlength\parskip{2.5mm}

\AtBeginDocument{
\addtolength{\abovedisplayskip}{2mm}
\addtolength{\abovedisplayshortskip}{2mm}
\addtolength{\belowdisplayskip}{2mm}
\addtolength{\belowdisplayshortskip}{2mm}
}

\color{MaterialGrey900}
\pagecolor{white}

\renewcommand{\familydefault}{\sfdefault}
\renewcommand{\baselinestretch}{1.5}

\makeatletter
\let\reftagform@=\tagform@
\def\tagform@#1{\maketag@@@{\ignorespaces\textcolor{MaterialGrey500}{(#1)}\unskip\@@italiccorr}}
\renewcommand{\eqref}[1]{\textup{\reftagform@{\ref{#1}}}}
\makeatother

\newcommand{\forceindent}{\leavevmode{\parindent=1em\indent}}

\usepackage{etoolbox}% http://ctan.org/pkg/etoolbox
\newtoggle{section}\togglefalse{section}% "In section" toggle
\newtoggle{atstarttoc}% "First entry in ToC" toggle
\newcounter{tocsubsection}% Keep track of first subsection
\makeatletter

\titlespacing\section{0mm}{12pt plus 2pt minus 2pt}{0pt plus 2pt minus 2pt}
\titlespacing\subsection{0pt}{12pt plus 2pt minus 2pt}{0pt plus 2pt minus 2pt}
\titlespacing\subsubsection{0pt}{12pt plus 2pt minus 2pt}{0pt plus 2pt minus 2pt}

\newcounter{subeq}
\renewcommand{\thesubeq}{\theequation\alph{subeq}}
\newcommand{\newsubeqblock}{\setcounter{subeq}{0}\refstepcounter{equation}}
\newcommand{\mysubeq}{\refstepcounter{subeq}\tag{\thesubeq}}

\title{Solving the One-Dimensional Wave Equation}
\author{Timo Vijn}
\date{September 21\textsuperscript{st} 2017 -- September 22\textsuperscript{nd} 2017}

\usepackage{hyperref}
\hypersetup{
    colorlinks,
    citecolor=MaterialGrey900,
    filecolor=MaterialGrey900,
    linkcolor=MaterialGrey900,
    urlcolor=MaterialGrey900
}

\usetikzlibrary{shapes.misc}

\usepackage[framed,numbered,autolinebreaks,useliterate]{mcode}

\usepackage{listings}

\usepackage{lscape}

\usepackage{fontspec}
% \setmainfont{[Fonts/MaisonNeue/MaisonNeue-Book.otf]}

\setmainfont [ Path = Fonts/MaisonNeue/ ] {MaisonNeue-Book.otf}
\setsansfont [ Path = Fonts/SF/ ] {SF-Compact-Text-Regular.otf}

% \setmainfont[
%  BoldFont={Minion Pro Bold}, 
%  ItalicFont={Minion Pro Italic},
%  BoldItalicFont={Minion Pro Bold Italic}
%  ]{Linux Libertine O}
 
 \setsansfont % load font from path
    [
        Path = Fonts/SF/,
        % Extension = .otf,
        UprightFont =       SF-Compact-Text-Regular.otf,
        BoldFont =          SF-Compact-Text-Medium.otf,
        ItalicFont =        SF-Compact-Text-RegularItalic.otf,
        BoldItalicFont =    SF-Compact-Text-MediumItalic.otf,
    ]
    {SF}

\begin{document}
\begin{titlepage}
    \begin{center}
        \vspace*{1cm}
        
        \textbf{Computational study of penetration mechanics of high performance\\Dyneema composites under ballistic loading}
        
        \vspace{0.5cm}
        A study of wave propagation and its dependence on model parameters
        
        \vspace{1.5cm}
        
        \textbf{T.P. (Timo) Vijn}
        
        \vfill
        
        A thesis presented for the degree of\\
        Bachelor of Science (BSc)
        
        \vspace{0.8cm}
        
        \includegraphics[width=0.3\textwidth]{tue}
        
        \vspace{0.6cm}
        
        Department of Mechanical Engineering\\
        \textit{Section of Mechanics of Materials}\\
        Eindhoven University of Technology\\
        Eindhoven, The Netherlands\\
        \textcolor{MaterialRedA400}{Version submitted for examination}\\
        February 4th 2018
        
    \end{center}
\end{titlepage}
\newpage
\section*{Preface}
In this bachelor’s thesis project the stress waves propagating as a result of ballistic impact through the body under impact have been studied. This project has been carried out between September 2017 and February 2018 at the Section of Mechanics of Materials of the Department of Mechanical Engineering at Eindhoven University of Technology.

The author would like to thank his examiners, R.H.J. (Ron) Peerlings, Associate Professor at the Department of Mechanical Engineering, and R. (Ramachandran) Varun Raj, Doctural Candidate (PhD) at the Department of Mechanical Engineering---not only for their help and support throughout this project, but also for showing him the wonders of their and their section's studies, which sparked his interest along the way.
\newpage
\setcounter{tocdepth}{2}
\tableofcontents
\newpage
\section{Introduction}
At the highest level, there are two motives that cover most of the research in the field of electricity price forecasting: the motive to improve the accuracy of EPF, and the motive to improve understanding of the (interplay of) factors that affect the accuracy of EPF.

As mentioned in [Catalão (2012)](https://www.notion.so/Catal-o-2012-2581f39d5513469a97ea001709f174f2), the price of electricity comes about as a consequence of many factors including its past values as well as exogenous variables such as load demand and production. The interest in such a quantity, that has numerous complicated and even non-observable factors at its core leads to a research space that is vast, and many factors that have been suggested to affect the electricity price are not fully established yet, let alone their complex interplay. Still, attention has steadily increased in recent years and research has proven that it is certainly possible to measure and model the effect of some of these factors to a greater or lesser extent. [Goodarzi et al. (2019)](https://www.notion.so/Goodarzi-et-al-2019-db841051dd2a4beea61e27e07f786951), for instance, seeks to model how forecast errors of renewable energy sources (RES) including wind and solar affect the imbalance volume and EPEX Spot price while taking no special interest in developing superior forecasting models. [Hagfors et al. (2016)](https://www.notion.so/Hagfors-et-al-2016-218dd1d0c9964dd59c3024205379441d) models the electricity price as a function of fundamental factors being the price of gas and coal and the carbon emission price, and finds that the sensitivities to these factors vary across the price distribution. Some factors seem to be very difficult to estimate however, for instance that of *market power* as pointed out in [Kolberg \& Waage (2018)](https://www.notion.so/Kolberg-Waage-2018-1bfb58f2df194567a53a3c9d77fedb0f). Although the behavior of participants is motivated partly by factors that are observable, factors remain that are unobservable, e.g. marginal production cost of individual participants, which introduce uncertainty.

On top of being subject to a multitude of factors, EPF is considerably region-dependent. Significant variation can be observed between electricity markets of different regions, even between those where similar regimes are in place, given inconsistent factors such as renewable energy policies, climate conditions, and economic environments. Many studies are therefore demarcated to a single region. Take [Monteiro et al. (2016)](https://www.notion.so/Monteiro-et-al-2016-565570927188470994eef1d584f88887) for instance, on the Iberian electricity market, [Hagfors et al. (2016)](https://www.notion.so/Hagfors-et-al-2016-218dd1d0c9964dd59c3024205379441d) on the UK market, [Kolberg \& Waage (2018)](https://www.notion.so/Kolberg-Waage-2018-1bfb58f2df194567a53a3c9d77fedb0f) on the Nordic market, and [Kath (2019)](https://www.notion.so/Kath-2019-72867b96bdcc4633ad33b9289996cdb2) and [Uniejewski et al. (2018)](https://www.notion.so/Uniejewski-et-al-2018-271aff7b89cd40deaa2461da8a1998d8) on the German market. What is more, [Cheng et al. (2016)](https://www.notion.so/Cheng-et-al-2016-aa38815d91e64d21b503c565e7045a54) and [Cheng et al. (2018)](https://www.notion.so/Cheng-et-al-2018-cb9ac9d5579144098ddb8e21ad53d0b4), on the market of the Chinese province of Yunnan, demonstrate that market regimes are not necessarily enforced on a nationwide level; market policies introduced in that province act as a pilot, from which valuable lessons can be drawn for later policy reform for the rest of China.

\section{Introduction}

































\end{document}